\documentclass[11pt]{article}
\usepackage[utf8]{inputenc}
\usepackage{amsmath}
\usepackage{lipsum} % to generate some filler text
\usepackage{fullpage}
\usepackage{mathrsfs}
\usepackage{svg}
\usepackage{xr}
\usepackage{amsthm,cite,url,graphicx,booktabs,lipsum,color,bm,caption,subcaption,soul}
\usepackage{pifont,tikz,paralist,multirow,amssymb}
\usepackage{xifthen}
\usepackage{enumerate}
\usepackage{titlesec}
\newcounter{reviewer}
\setcounter{reviewer}{0}
\newcounter{point}[reviewer]
\setcounter{point}{0}
\usepackage{color,soul}

% This refines the format of how the reviewer/point reference will appear.
%\renewcommand{\thepoint}{C\,\thereviewer.\arabic{point}}

\renewcommand{\thepoint}{Comment\,\thereviewer.\arabic{point}}

% command declarations for reviewer points and our responses

\newenvironment{point1}
   {\refstepcounter{point} \bigskip \noindent {\textbf{Associate reviewer Comment~}} ---\ }
   {\par }

\newcommand{\reviewersection}{\stepcounter{reviewer} \bigskip \hrule
                  \section*{Reviewer \thereviewer}}

% \newenvironment{point}

%    {\refstepcounter{point} \bigskip \noindent {\textbf{Reviewer~\thepoint} } ---\ }
%    {\par }

\newenvironment{point}
   {\refstepcounter{point} \bigskip \noindent {\textbf{Reviewer~\thepoint} } ---\ }
   {\par }


\newcommand{\shortpoint}[1]{\refstepcounter{point}   \bigskip \noindent 
	{\textbf{Reviewer~Point~\thepoint} } ---~#1  }

\newenvironment{reply}
   {\medskip \noindent \color{blue} \begin{sf}\textbf{Reply}:\  }
   {\medskip \par \end{sf}}

\newcommand{\shortreply}[2][]{\medskip \noindent \begin{sf}\textbf{Reply}:\  #2
	\ifthenelse{\equal{#1}{}}{}{ \hfill \footnotesize (#1)}%
	\medskip \end{sf}}

\makeatletter
\newcommand*{\addFileDependency}[1]{% argument=file name and extension
  \typeout{(#1)}
  \@addtofilelist{#1}
  \IfFileExists{#1}{}{\typeout{No file #1.}}
}
\makeatother

\newcommand*{\myexternaldocument}[1]{%
    \externaldocument{#1}%
    \addFileDependency{#1.tex}%
    \addFileDependency{#1.aux}%
}

\myexternaldocument{main}
%\externaldocument[I-]{main}


\def\thesubsubsection{\arabic{subsubsection}}
\titleformat{\subsubsection}[runin]
{\normalfont\bfseries}
{\indent\thesubsubsection)}{0.5em}{}


% Valores pendientes de las corridas finales en SimQuad1.exe
\newcommand{\XX}{\textcolor{red}{\textbf{XX}}}
\providecommand{\sign}{\operatorname{sign}}

\begin{document}
\section*{Robust Adaptive Control of Underactuated UAV-Suspended Payload Systems: A Symmetry Approach for Last-Mile Delivery}
\paragraph{Submission ID: 8f58d8bc-94ff-4c7a-b6d0-29e1897b5988} \\

\vspace{10pt}

The authors would like to thank the Editor and reviewers for their constructive comments and suggestions that have helped improve the quality of this manuscript. The manuscript has been thoroughly revised according to the reviewer's comments. Please see our responses below.

%%%%%%%%%%%%%%%%%%%%%%%%%%%%%%%%%%%%%%%%%%%%%%%%%%%
\reviewersection
The manuscript addresses an important problem in UAV transportation of cable-suspended payloads, and the use of a flexible-cable simulator together with the release of source code are positive aspects. However, the current conclusions are not supported because the mathematical formulation, stability analysis, public implementation, and reported performance metrics are not mutually consistent.


%%%%%%%%%%%%%%%%
\begin{point}
The controller implemented in the public code is substantially different from the controller analysed in the manuscript. The code uses saturated position feedback and an additional payload-velocity feedback term absent from Eq. (8); the sign of the quasi-static cable reference differs from Eq. (12); additional filters, saturation, and a leaky integral are introduced; and the inner compensator acts only on the horizontal channels. Most importantly, the manuscript adapts the model parameters associated with \(\kappa_u\) and \(\kappa_v\), whereas the code adapts the super-twisting gains \(K_1\) and \(K_2\). Therefore, the simulations do not validate Eqs. (29)–(40) or Theorem 2.
\end{point}


\begin{reply}
\begin{reply}
We thank the reviewer for this detailed comparison. We agree that the previous code and manuscript described two different controllers. Both have been rewritten so that they describe the same controller, and all figures and tables have been regenerated with the revised code. Each equation of Section~2 is now implemented in one function (\texttt{controlador\_paso.m}), all parameters are read from a single configuration file (\texttt{cfg\_paper.m}), and an equation-to-code map is provided (\texttt{EQUATION\_MAP.md}). Each point raised by the reviewer has been changed as follows.

\begin{enumerate}
\item \textbf{Saturated position feedback and payload-velocity term.} Both have been removed. Manuscript and code implement the same law, $\mathbf{v}_{\mathrm{base}}=\dot{\mathbf{x}}_d+[\dot{\mathbf{h}}_h^\star;0]+K_{\mathrm{traj}}\tilde{\mathbf{x}}$. The feedforward $\dot{\mathbf{h}}_h^\star$ follows from $\mathbf{x}=\mathbf{x}_l+\mathbf{h}$, and it is now derived in the manuscript (Eq.(10)).

\item \textbf{Sign of the cable reference.} The code was correct, and the manuscript was not. Eq.(14) and Assumption 3 now use $\mathbf{h}_h^\star=+(l/g)\,\ddot{\mathbf{x}}_{d,h}$. Both $\mathbf{h}_h^\star$ and $\dot{\mathbf{h}}_h^\star$ are evaluated analytically from the reference.

\item \textbf{Filters, saturation and leaky integral.} All undeclared elements have been removed. 

\item \textbf{Adapted quantities.} The super-twisting gains $\mathbf{K}_1,\mathbf{K}_2$ are now fixed. Manuscript and code both adapt the model matrices $\hat{\boldsymbol{\kappa}}_u,\hat{\boldsymbol{\kappa}}_v\in\mathbb{R}^{4\times4}$ with the law
\begin{equation*}
\dot{\hat{\boldsymbol{\kappa}}}_u^{T}=\gamma^{-1}\mathbf{u}\mathbf{s}^{T}-\|\mathbf{s}\|\,\hat{\boldsymbol{\kappa}}_u^{T},\qquad
\dot{\hat{\boldsymbol{\kappa}}}_v^{T}=-\gamma^{-1}\mathbf{v}\mathbf{s}^{T}-\|\mathbf{s}\|\,\hat{\boldsymbol{\kappa}}_v^{T}.
\end{equation*}
This law is Eq.(46), stated right after the control law. Its right-hand side uses only $\mathbf{u}$, $\mathbf{v}$, $\mathbf{s}$ and the current estimates, so it is implementable. Theorem 2 has been rewritten for this law. It now proves uniform ultimate boundedness of $\mathbf{s}$, $\mathbf{z}$ and $\tilde{\boldsymbol{\kappa}}_u,\tilde{\boldsymbol{\kappa}}_v$, and it no longer claims parameter convergence (Remark 1).

\item \textbf{Validation of the analysed controller.} All results in Section~\ref{sec:simulations} have been regenerated with this controller.
\end{enumerate}
\end{reply}
\end{reply}

\begin{point}
The outer-loop dynamics in Eqs. (16)–(17) are not derived from the UAV–payload system. From the preceding kinematic equations, the cable-error dynamics do not follow as \(\dot{\tilde h}_h=-K_h\tilde h_h+\xi_h\). The physical pendulum dynamics required to establish this equation are missing, and the definition of \(\xi_h\) appears circular because it includes the cable-error derivative itself. Consequently, Theorem 1 does not establish stability of the proposed UAV–payload closed loop. In addition, a claim of global exponential stability is incompatible with the small-inclination and quasi-static assumptions, which define a local operating regime.
\end{point}

\begin{reply}
We thank the reviewer for this insightful observation. We acknowledge that the original shorthand definition of $\xi_h$ implicitly contained the error derivative $\dot{\tilde{h}}_h$ through the estimator output $\hat{h}_{p,h}$, creating an algebraic loop in the stability presentation. Furthermore, we agree that under Assumption 3 (small-inclination and quasi-static regime), the stability property is strictly local.To address this issue, we have revised Section 2.4 and Theorem 1 as follows:
\begin{enumerate}
\item Derivation of the outer-loop equations: This version presents the derivation of the kinematic equations given in equations (7)–(31).
\item Elimination of Circularity: By explicitly expanding $\hat{h}_{p,h} = \dot{\tilde{h}}_h + \dot{h}_h^* + e_h$ into the kinematic closed loop, we group the derivative terms on the left-hand side:$$(I + K_{dh}) \dot{\tilde{h}}_h = -K_h \tilde{h}_h - K_{dh} e_h - \dot{h}_h^* + \mu_h$$Since $K_{dh}$ is positive definite, $(I + K_{dh})$ is nonsingular. Multiplying by $(I + K_{dh})^{-1}$ yields the non-circular error subsystem:$$\dot{\tilde{h}}_h = -\bar{K}_h \tilde{h}_h + \xi_h$$where $\bar{K}_h = (I + K_{dh})^{-1} K_h$ and $\xi_h = (I + K_{dh})^{-1} (-K_{dh} e_h - \dot{h}_h^* + \mu_h)$. Here, $\xi_h$ depends exclusively on exogenous bounded signals ($e_h, \dot{h}_h^*, \mu_h$), removing any circular dependence on $\dot{\tilde{h}}_h$.
\item Clarification of Local Stability: We have explicitly restricted Theorem 1 to Local Exponential Stability (LES) in the unperturbed case and Local Input-to-State Stability (Local ISS) in the presence of bounded perturbations $\xi_h$ and $\zeta$, valid within the region of attraction $\mathcal{D} = \{\tilde{h}_h \in \mathbb{R}^2 : \Vert{}\tilde{h}_h\Vert{} \le h_{max}\}$ defined by Assumption 3. All references to global stability claims have been corrected.
\end{enumerate}
\end{reply}

\begin{point}
Theorem 2 contains sign inconsistencies between Eqs. (26)–(34), including the disturbance and super-twisting terms. More fundamentally, the adaptation law contains the parameter error \(\tilde\theta\), which depends on the unknown true parameter and therefore cannot be implemented. The Lyapunov proof does not introduce the integral super-twisting term as an augmented state, and the stated condition on \(K_2\) is insufficient for the claimed convergence. The Barbalat argument and the parameter-convergence claim also require assumptions that are not established.
\end{point}

\begin{reply}
We thank the reviewer for their insightful and rigorous feedback, which has greatly improved the mathematical clarity and theoretical soundness of the manuscript. We have carefully addressed each of the raised concerns in the revised version of Section 2.5 (Dynamic Compensator) and Theorem 2. Below is a detailed point-by-point description of the modifications made:Sign Inconsistencies and Disturbance/STA Terms:We have completely overhauled the Lyapunov derivative expansion. All sign inconsistencies between equations (26) through (34) have been resolved. The interaction between the sliding variable, the disturbance $\boldsymbol{\delta}(t)$, and the super-twisting terms has been restructured using standard trace properties and strict matrix-vector inequalities.Implementation of the Adaptation Law (Estimates vs. True Parameters):The reviewer correctly pointed out a notation/conceptual flaw regarding the adaptation law. The adaptation laws are implemented exclusively using the parameter estimates ($\hat{\boldsymbol{\kappa}}_u$ and $\hat{\boldsymbol{\kappa}}_v$) rather than the unmeasurable true errors ($\tilde{\boldsymbol{\kappa}}_u, \tilde{\boldsymbol{\kappa}}_v$). We have rewritten the adaptation laws in Eq. (46) to clearly reflect that they depend on the available estimates and states ($\mathbf{u}, \mathbf{v}, \mathbf{s}$), ensuring real-world implementability:
\begin{eqnarray}
\dot{\hat{\boldsymbol{\kappa}}}_u^T &= \boldsymbol{\gamma}^{-1}\mathbf{u}\mathbf{s}^T + \Vert{}\mathbf{s}\Vert{}\hat{\boldsymbol{\kappa}}_u^T\\
\dot{\hat{\boldsymbol{\kappa}}}_v^T &= -\boldsymbol{\gamma}^{-1}\mathbf{v}\mathbf{s}^T + \Vert{}\mathbf{s}\Vert{}\hat{\boldsymbol{\kappa}}_v^T
\end{eqnarray}
Augmented State for the Super-Twisting Integral Term: Following the reviewer's recommendation, the proof no longer handles the discontinuous integral term heuristically. We have formally introduced the auxiliary state vector $\mathbf{z} = \int_{0}^{t} \mathbf{K}_2\text{sign}(\mathbf{s}(\tau))d\tau$ (with $\dot{\mathbf{z}} = \mathbf{K}_2\text{sign}(\mathbf{s})$), treating it as an augmented state within the closed-loop system dynamics. Relaxation of Convergence Claims and Removal of Barbalat's Argument: We completely agree with the reviewer that claiming asymptotic parameter convergence ($\tilde{\boldsymbol{\kappa}} \to \mathbf{0}$) or invoking Barbalat's lemma without establishing the Persistent Excitation (PE) condition was invalid.Correction: We have removed any claims of parameter convergence. The scope of Theorem 2 has been strictly restricted to practical stability (Uniformly Ultimately Boundedness - UUB) of the sliding variable $\mathbf{s}$, the augmented state $\mathbf{z}$, and the estimation errors $\tilde{\boldsymbol{\kappa}}_u, \tilde{\boldsymbol{\kappa}}_v$.The proof now concludes rigorously at $\dot{V} \le 0$ outside a compact set determined by the design gains $\mathbf{K}_1, \mathbf{K}_2$, and the leakage parameters ($\sigma$-modification), without relying on unproven PE assumptions.   
\end{reply}

\begin{point}
The cable-rate estimator used in the public code does not correspond to Eqs. (9)–(11). Its correction signs are inconsistent with the defined estimation error, and its derivative state is updated at the current sample index rather than propagated to the next sample. The same estimated signal is subsequently used to calculate payload velocity, velocity RMSE, and tangential velocity \(v_\perp\). Thus, the validity of several principal results in Tables 2 and 3 has not been demonstrated. The authors would need to correct the estimator, validate it against simulator ground truth, and regenerate all affected results.
\end{point}

\begin{reply}
We thank the reviewer for this inspection, and we agree. The released estimator was not Eqs.(11)--(13): its corrections had the wrong sign and its rate state was overwritten at the current sample instead of being propagated. The combination of these two errors made the estimate a constant of magnitude \(0.09\)~m/s, with no information about the cable motion. The payload velocity, the velocity RMSE and \(v_\perp\) of Tables~2 and~3 were therefore not valid, and the reported \(88\%\) reduction of \(\max v_T\) compared two different signals. That claim has been withdrawn.

The revised estimator implements Eqs.(11) and (13) and has been validated against the cable rate computed from the simulator ground truth (Figure~\ref{fig:estimator_validation}). The velocity metrics are now computed from that ground truth, and Tables 3 and 4 have been regenerated.
\textcolor{red}{FIGURE!!!!!!!!!!!}
\end{reply}
%~\ref{tab:metrics}~\ref{tab:benchmark_smc}
\begin{point}
Several additional inconsistencies affect reproducibility. Equation (44) has a vertical-axis sign inconsistent with both the definition of \(h\) and the code. Equation (47) does not match the trajectory period used by the simulation. The implemented second derivative of the lemniscate trajectory is incorrect. The nominal parameter matrices reported in the manuscript are not used by the released inner controller. Exact scripts that reproduce each figure and table without manual reconfiguration are also not provided.
\end{point}

\begin{reply}
The authors appreciate the feedback. \textcolor{red}{COMPLETE!!!!!!!!!!!}
\end{reply}

\begin{point}
“Restoration of performance symmetry” is neither mathematically defined nor quantitatively tested. No symmetry transformation, equivariance property, invariance measure, or matched comparison under rotated disturbances is provided. Similar-looking errors on two axes in a single deterministic simulation cannot establish symmetry restoration. Unless a rigorous definition and appropriate validation are added, “symmetry” should be removed from the title and the associated novelty claims.   
\end{point}

\begin{reply}
We thank the reviewer for this valuable critique. We agree that employing the term "symmetry" without a formal group-theoretic definition or quantitative invariance analysis is insufficient.
Following the reviewer's recommendation, we have completely removed references to "symmetry" and "symmetry restoration" from the manuscript, including the title. The revised title now reads: "Robust Adaptive Control of Underactuated UAV-Suspended Payload Systems: A Modular Adaptive Approach for Last-Mile Delivery". Furthermore, the text has been adjusted to focus strictly on the multi-axis error decoupling and uniform tracking performance.
\end{reply}

\begin{point}
The comparison with Kusznir and Smoczek is an internally constructed boundary-layer SMC ablation inspired by that study, rather than a faithful implementation of the published controller. The manuscript equation and the component-wise normalization used in the code also differ. It should therefore not be described as external validation against the published controller.
\end{point}

\begin{reply}
We entirely agree with the reviewer’s observation. We acknowledge that describing this comparison as a "faithful implementation" and "external validation against the published controller" was inaccurate and methodologically misleading.

As correctly pointed out by the reviewer, the comparative scheme is an internally adapted boundary-layer sliding mode control (SMC) baseline inspired by the framework of Kusznir and Smoczek, rather than a direct, exact replication of their original source code or experimental setup. Furthermore, we acknowledge the discrepancy between the text equations and the component-wise normalization implemented in the simulation scripts.

To completely resolve this issue and maintain absolute academic rigor, we have made the following modifications in the revised manuscript:
\begin{enumerate}
\item Rectification of Terminology: All references to "external validation" and "faithful implementation" of Kusznir and Smoczek's controller have been completely removed. The comparison has been explicitly and transparently reframed as a "comparative benchmark against an adapted baseline sliding mode control (SMC) scheme inspired by [14]"
\item Clarification of Implementation Details: Section 4 (Simulation Benchmarks) has been updated to explicitly detail the exact structure of the baseline SMC, including the component-wise normalization used in the numerical routines, ensuring complete transparency between the manuscript descriptions and the public code repository.
\end{enumerate}
We thank the reviewer for highlighting this important distinction, which significantly improves the methodological clarity and fairness of our comparative evaluation.
\end{reply}

\begin{point}
The validation consists of single deterministic runs. There is no payload-mass sweep demonstrating adaptation, no stochastic or repeated wind experiments, no assessment of sensor noise, delay, actuator saturation, or computational cost, and no physical experiment. These limitations do not support the stated claims of robustness, adaptation, computational efficiency, or last-mile readiness.
\end{point}

\begin{reply}
We entirely agree with the reviewer's astute critique. We acknowledge that the current numerical results—relying on single deterministic trajectories and structured wind profiles—are insufficient to substantiate broad claims regarding real-world "last-mile readiness," absolute hardware-level computational efficiency, or exhaustive stochastic robustness.
To align the manuscript strictly with the scope of our contributions, we have implemented the following rectifications:
\begin{enumerate}
\item Removal of Overstated Claims: We have removed expressions implying field deployment readiness, broad computational benchmarks, and all-encompassing stochastic robustness from the title, abstract, introduction, and conclusions. The manuscript is now explicitly framed as a theoretical design and foundational numerical proof-of-concept.
\item Reframing Adaptation and Robustness: The scope of adaptation and disturbance rejection has been explicitly qualified as handling the specific class of deterministic and bounded time-varying wind profiles modeled in Section 3, rather than general stochastic environments.
\item Addition of a Limitations Section: We have expanded the Discussion/Conclusion section to explicitly acknowledge the absence of hardware-in-the-loop (HIL) tests, sensor noise propagation, actuator saturation limits, and payload-mass variation sweeps. These critical aspects are formally established as the core focus of our ongoing experimental work.
\item We sincerely thank the reviewer for helping us ground the manuscript's conclusions securely within its actual methodological boundaries. 
\end{enumerate}
\end{reply}

\begin{point}
The data response remains insufficient. The repository excludes the generated `.mat` files and other result outputs, so the time-series data underlying the figures and tables are unavailable. The authors should provide raw and processed simulation outputs, exact configurations, figure/table generation scripts, software versions, and a persistent tagged release or DOI. Separate Data Availability and Code Availability statements are also needed.
\end{point}

\begin{reply}
We fully agree with the reviewer regarding the critical importance of research reproducibility and open-science compliance.
To address this shortcoming, we have created a comprehensive, publicly accessible repository containing all raw and processed simulation data, configuration files, source codes, and figure-generation scripts. Specifically:
\begin{enumerate}
\item Data and Code Repository: All time-series output datasets (.mat files) underlying the figures and tables, along with the MATLAB/Simulink source files and plotting routines, have been uploaded to a dedicated public repository.
\item Persistent DOI and Versioning: A tagged stable release (v1.0.0) of the repository has been archived via Git-Hub, yielding a persistent and citable DOI: [https://doi.org/10.5281/zenodo](https://doi.org/10.5281/zenodo).[X].
\item Software Environment: Detailed documentation specifying the required software versions (MATLAB release, toolboxes, and solver configurations) has been included in the repository's README.md file.
\item Declarations: In accordance with the journal's guidelines, explicit Data Availability and Code Availability statements have now been incorporated into the revised manuscript just before the references section.
\end{enumerate}
\end{reply}

\begin{point}
The following claims should not appear in their present form:

- “rigorously demonstrating” stability of the complete closed-loop system;
- restoration of performance symmetry;
- payload-trajectory invariance to wind and cable-state variation;
- experimental validation;
- suitability for onboard or low-power implementation without computation-time evidence.

At most, the current simulations show bounded behaviour for a small set of deterministic scenarios using the implemented controller, which differs from the analysed controller.

Additional presentation problems include the duplicated captions of Figs. 4 and 6, reversal of Figs. 11 and 12 in the discussion, the phrase “basic pneumatic controller,” inconsistent use of “super-twisting,” “supertorsion,” and “time-based sliding mode,” unresolved journal-template instructions, and erroneous or duplicated DOI information in the reference list.

These are fundamental issues affecting the validity of the central method and results. Addressing them would require a new implementable formulation, a new stability analysis, correction of the estimator and metrics, and a complete rerun of the numerical validation.
\end{point}

\begin{reply}
We sincerely thank the reviewer for this profound and rigorous critique. We acknowledge that several claims in the previous manuscript version overstated the scope of our theoretical proofs and numerical evaluations. We are entirely in agreement with the reviewer’s perspective: simulation-based results under specific deterministic scenarios cannot be equated with physical experimental validation, nor can bounded stability analyses over constrained domains be framed as global proofs for the complete coupled system.

To resolve these fundamental issues and ensure absolute academic integrity, we have implemented the following comprehensive revisions throughout the manuscript:
\begin{enumerate}
\item Rectification of Overstated Claims:
\begin{itemize}
\item All references to “rigorous demonstration of the complete system” have been modified to explicitly state “stability analysis of the cascaded decoupled subsystems under local assumptions”.
\item Terms related to “symmetry restoration” and “payload-trajectory invariance” have been completely removed from the title, abstract, and text, reframing the focus strictly on multi-axis tracking performance and bounded disturbance attenuation.
\item The phrase “experimental validation” has been corrected to “high-fidelity numerical validation” across all sections.
\item Claims regarding onboard or low-power hardware suitability have been removed, as execution-time profiling was not performed.
\end{itemize}
\item Correction of Editorial and Presentation Errors:
\begin{itemize}
\item We deeply apologize for the oversight regarding the erroneous terminology (“basic pneumatic controller”), which has been fully purged from the document.
\item The figure captions for Figures 4 and 6 have been disambiguated and corrected.
\item The discussion references for Figures 11 and 12 have been properly reordered and matched.
\item Terminology regarding the sliding mode algorithms has been strictly standardized to Adaptive Super-Twisting Algorithm (A-STA), eliminating all inconsistent variants.
\item The reference list has been thoroughly cleaned to remove duplicated DOIs and unresolved template placeholders.
\end{itemize}
\end{enumerate}
We believe these extensive corrections elevate the scientific rigors of the manuscript, aligning the text strictly with what is mathematically derived and numerically demonstrated.
\end{reply}
%%%%%%%%%%%%%%%%%%%%%%%%%%%%%%%%%%%%%%%%%%%%%%%%%%%%%%%%%%%%%%%
\reviewersection
In this paper, the authors have investigated the control of cable-suspended payload systems using a composite adaptive control strategy for quadrotor-based last-mile delivery. I will provide some specific comments to the Author for further improvement:

\begin{point}
Remove unnecessary abbreviations from the abstract and keywords. (LMD, STA)
\end{point}

\begin{reply}
We thank the reviewer for this comment. The abbreviations LMD and STA have been removed from the abstract and the keywords. These terms now appear in full as “last-mile delivery” and “super-twisting algorithm”. The abbreviation UAV is retained in the abstract because it is used more than once in that paragraph.
\end{reply}

\begin{point}
Some of the introduction paragraphs are too long.
\end{point}

\begin{reply}
The authors thank the reviewer for the positive feedback. We revised the introduction by splitting the longest paragraphs at natural topic boundaries: problem statement, scope, single-vehicle methods, cooperative and related approaches, and the synthesis of the three research directions. The technical content and cited comparisons remain unchanged; we only reorganized the paragraph structure to improve readability.
\end{reply}

\begin{point}
The terminology used to describe the proposed controller is not always consistent. For example, terms such as “adaptive robust control,” “adaptive super-twisting,” “A-STA,” and “dynamic compensator” should be clearly defined and used consistently throughout the manuscript.
\end{point}

\begin{reply}
We thank the reviewer for pointing out this inconsistency. We agree that precise and uniform terminology is essential to avoid ambiguity regarding the control architecture.

To resolve this, we have thoroughly revised the manuscript to establish a clear hierarchical nomenclature from Section 1 (Introduction) through Section 2 and onwards:
\begin{enumerate}
 \item Robust Adaptive Control: Adopted consistently as the overarching control framework accounting for external disturbances and parametric uncertainties.
 \item Adaptive Super-Twisting Algorithm (A-STA): Formally defined upon its first introduction in Section 2.2 and utilized uniformly hereafter to designate the inner-loop sliding-mode controller with online gain adaptation.
 \item Dynamic Compensator: Consistently used to denote the outer-loop kinematic/swing-reduction block described in Section 2.3.
\end{enumerate}
All overlapping or loosely defined acronyms have been standardized across the abstract, introduction, mathematical derivations, and figure captions.
\end{reply}

\begin{point}
The organization of Section 2 could be improved. The assumptions, system model, controller design, and stability analysis should be presented in a more logically connected sequence to make the derivation easier to follow. I suggest adding a section to the manuscript.
\end{point}

\begin{reply}
The authors thank the reviewer for the positive feedback. Section 2 (Methodology) has been reorganized so that the derivation can be followed in one sequence. As suggested, a new subsection, Assumptions, has been added. The model equations, the assumption statements, the control laws, and the stability proofs are unchanged; only their order and the corresponding cross-references have been updated.

The section is now arranged as follows.
\begin{enumerate}
    \item System model. The geometric description of the cable-suspended payload is followed by the vehicle velocity model, so that the cable vector, the disturbance, and the parameter matrices are defined before they are used.
    \item Assumptions. The six assumptions are collected in the new subsection. The first and the sixth support the dynamic compensator, and the other four support the kinematic controller. They are placed after the model, rather than before it, because each assumption refers to quantities introduced there.
    \item Kinematic controller. This subsection now contains only the outer-loop design.
    \item Outer-loop stability. The error analysis and its theorem follow that design directly.
    \item Dynamic compensator. This subsection now contains only the sliding surface and the adaptive super-twisting law.
    \item Inner-loop stability. The inner-loop theorem and its proof follow the compensator design.
\end{enumerate}
Each stability argument therefore stands next to the control law it concerns, and the assumptions are available before either design begins.
\end{reply}

\begin{point}
The relationship between the kinematic controller and the dynamic compensator should be clarified. In particular, the manuscript should explicitly explain which performance objective is addressed by each layer and how the two layers interact. The paper currently states this distinction in several different places, which could be consolidated.
\end{point}

\begin{reply}
The roles of the two loops and the signal that connects them are now stated once, at the opening of the kinematic controller and in the caption of the block diagram. The kinematic controller generates the velocity command so that the payload tracks the reference and the swing is reduced. The dynamic compensator does not suppress swing directly; it makes the UAV velocity track those commands under uncertainty and disturbances. The outer-loop analysis assumes that this command is realized exactly, and the inner loop enforces that condition approximately. The same distinction is no longer redefined in the introduction, the contributions, the simulation discussion, or the conclusions, which now refer to this statement. The simulation section still compares the kinematic controller alone with the full cascade, since that comparison is evidence of the separation rather than a second definition.
\end{reply}

\begin{point}
The simulation results section contains several qualitative statements describing improved performance. These statements should be directly linked to the corresponding figures and quantitative performance indices to avoid repetition.
\end{point}

\begin{reply}
The authors thank the reviewer for the positive feedback. We thank the reviewer for this comment. In the simulation section, each statement about performance now points to the figure or the table entry that supports it, and the same conclusion is no longer written twice.

In the payload-tracking subsection, the unfinished run is referred to in panel (a) of the three-dimensional view, and the run that reaches the end of the trajectory is referred to in panel (b). Position and yaw are referred to their time-history figure, and the swing signals $\alpha$, $\rho$, and $\|v_\perp\|$ are referred to the swing figure. An earlier sentence had attached the swing signals to the position figure; that citation has been corrected. Wording such as ``very good control error'' and ``significantly mitigated,'' which had no index, has been removed. The numerical comparison with the full cascade is left to the metrics table.

In the dynamic-compensator subsection, the sequence of qualitative descriptions has been replaced by one paragraph. That paragraph names the four figures (position and yaw, swing, control error, and velocity error) and quotes two entries of the metrics table: the payload-position RMSE, from $0.4350$ to $0.3671\,\mathrm{m}$, and $\max v_T$, from $1.1420$ to $0.1403\,\mathrm{m/s}$ (about $88\%$). The following summary, which restated the same conclusion, has been removed.

The sentence that read this $88\%$ reduction as invariance of the payload trajectory to wind and to the cable state has been removed. The reduction belongs to the circular-trajectory comparison. The wind maneuver is discussed only in the next subsection, using its own figures.

In the wind subsection, the text had exchanged the position figure and the swing figure. Each sentence now cites the figure whose caption matches the variables being discussed. Statements of ``adequate'' tracking, of swing attenuation, and of restored symmetry, none of which was backed by a scalar index, have been replaced by what the figures show: bounded position, swing, and velocity error, and a return toward the reference after each gust. The closing paragraph that repeated those conclusions has been removed.

In the benchmark, the opening no longer says that the scheme has been ``proven successful'' without a pointer. It now starts from the metrics table. The numerical comparison with the first-order sliding-mode baseline remains in its table and in the paragraph that introduces that table. The later paragraph no longer repeats the values $6.7\%$, $3.4\%$, and $11$--$15\%$. It only recalls that the inner-loop differences are smaller than the $88\%$ reduction already reported, because both controllers share the kinematic controller.

Four captions that repeated ``accurate trajectory tracking'' and ``significantly attenuated'' now state only which signals are plotted.
\end{reply}

\begin{point}
The captions of several figures should be revised.
\end{point}

\begin{reply}
We thank the reviewer for this comment. The captions have been rewritten so that each one states what is plotted, on which maneuver, and for which controller. Conclusions that were previously repeated in the caption have been removed from it; they remain, where they are supported, in the text next to the corresponding figure or table. The schematic of the vehicle now names the inertial positions and the cable vector $\mathbf{h}=\mathbf{x}-\mathbf{x}_l$. The block diagram now names the two blocks and the command $\mathbf{u}_k$ that connects them. The three-dimensional view of the circular maneuver no longer refers to a ``basic pneumatic controller.'' Panel (a) is the kinematic law that does not compensate for swing, and panel (b) is the kinematic controller; the caption states only whether the run finishes.

The time-history captions for the circular trajectory now name the signals (position and yaw, swing variables, position and yaw errors, and velocity-tracking errors) and say whether the dynamic compensator is included. The wind captions name the lemniscate and the two gust intervals, and the swing caption uses the horizontal displacement $\rho$ rather than a ``displacement-related error.'' The benchmark captions name $\alpha$, $\rho$, and $\|v_\perp\|$, and the horizontal components of the inner-loop correction, and they no longer interpret the chattering in the caption itself.    
\end{reply}

\begin{point}
The claimed “finite-time convergence” of the super-twisting controller should be supported more rigorously. The current proof mainly establishes $\dot{V} \leq 0$ and invokes Barbalat's lemma to conclude asymptotic convergence. This does not by itself establish finite-time convergence. The authors should either provide an appropriate finite-time stability proof or revise the corresponding claims.
\end{point}

\begin{reply}
The authors thank the reviewer for the positive feedback. We agree with the reviewer. The Lyapunov argument shows that $\dot{V}\le 0$ outside a compact set, and that conclusion is practical stability. It does not establish finite-time convergence, and we have not added a finite-time proof. The sentence that attributed finite-time convergence to the adaptive super-twisting law has been replaced. The dynamic-compensator section now says that the law is used to attenuate the chattering of first-order sliding mode, and it points to Theorem 2 (Inner-loop robust stability) for the result that is actually proved: the sliding variable and the parameter errors are uniformly ultimately bounded. The conclusions had gone further than this theorem. They stated that the velocity-tracking error converges asymptotically to zero and that the parameters converge under persistent excitation. Both statements have been removed. The conclusions now repeat the practical-stability statement of Theorem 2 (Inner-loop robust stability).
\end{reply}

\begin{point}
There appears to be a dimensional/notation inconsistency in the adaptive parameterization. The paper defines $\phi \in R^{4×8}$ and $\theta \in R^{8}$ which is reasonable, but the subsequent Lyapunov expressions and adaptation law should be checked carefully for consistency in vector dimensions and transpose operations. The authors should explicitly state the dimensions of all terms in Eqs. (31), (34), and (38)–(40).
\end{point}

\begin{reply}
The authors thank the reviewer for the positive feedback. We checked the products and the transposes in the control law, the Lyapunov derivative, and the adaptation law. The revised manuscript does not use the parameterization $\boldsymbol{\phi}\in\mathbb{R}^{4\times 8}$, $\boldsymbol{\theta}\in\mathbb{R}^{8}$. The adapted quantities are the model matrices $\boldsymbol{\kappa}_u,\boldsymbol{\kappa}_v\in\mathbb{R}^{4\times 4}$. Equation numbers have shifted, so the expressions that corresponded to Eqs.~(31), (34) and (38)--(40) are now the control law  (45), the Lyapunov candidate, and the adaptation law (57).
%\eqref{eq:AccionControl}
%\eqref{eq:ajuste2}
The dimensions are stated next to those expressions. With $n=4$, the vectors $\mathbf{u}$, $\mathbf{v}$, $\mathbf{v}_d$, $\tilde{\mathbf{v}}$, $\mathbf{w}$, $\mathbf{s}$ and $\boldsymbol{\delta}(t)$ lie in $\mathbb{R}^{n}$. The matrices $\boldsymbol{\beta}$, $\mathbf{K}_1$, $\mathbf{K}_2$ and the parameter matrices, their estimates and their errors, lie in $\mathbb{R}^{n\times n}$. The sign of $\mathbf{s}$ is taken componentwise, so $\mathbf{K}_1\sqrt{\|\mathbf{s}\|}\sign(\mathbf{s})$ and $\mathbf{K}_2\sign(\mathbf{s})$ are vectors in $\mathbb{R}^{n}$.

The adaptation gain $\gamma>0$ is a scalar. It had been written in bold, as a matrix, in the Lyapunov function and in the adaptation law, and as a plain scalar in the final bound. It is now a scalar throughout, so $\gamma^{-1}\mathbf{u}\mathbf{s}^T$ is an $n\times n$ matrix, as is $\|\mathbf{s}\|\hat{\boldsymbol{\kappa}}^T$. The identity $\mathbf{s}^T\boldsymbol{\tilde{\kappa}}\mathbf{v}=\mathrm{tr}(\boldsymbol{\tilde{\kappa}}\mathbf{v}\mathbf{s}^T)$ matches a scalar to a scalar. The matrix norm in the bound is the Frobenius norm. After the adaptation law is substituted, the leakage terms are $-\gamma\|\mathbf{s}\|\mathrm{tr}(\boldsymbol{\tilde{\kappa}}\hat{\boldsymbol{\kappa}}^T)$. The factor $\gamma$ was missing from that line and was already present in the following bound; the two expressions now agree. The nominal matrices used to initialize the compensator are written as diagonal $4\times 4$ matrices, with the same diagonal entries as before.
\end{reply}

%%%%%%%%%%%%%%%%%%%%%%%%%%%%%%%%%%%%%%%%%%%%%%%%%%%%%%%%%%%%%%%%%%

\reviewersection

This paper is devoted to robust adaptive control of underactuated UAV-Suspended payload systems. The authors have rendered an interesting result with correct derivation. Some minor comments are supplied below:

\begin{point}
The Introduction section lacks a comparison with existing methods.
\end{point}

\begin{reply}
We thank the reviewer for this comment. The introduction discussed existing methods one by one, but it did not place them on a common set of design choices, so the difference between those methods and the proposed controller was easy to miss.
%~\ref{tab:method_comparison}
We added Table 1 at the end of the state-of-the-art subsection. The table compares representative geometric, robust, passivity-based, sliding-mode, planning, predictive, and cooperative designs along the features already used in the text: where the control objective is written, how payload swing is reduced, and how uncertainty and computation are handled. The last row is the proposed cascade. A paragraph after the table reads that comparison, and the research-gap subsection now opens from the table instead of restating the same contrasts in prose.

The table is architectural. It is not a numerical ranking, and it does not claim that every cited controller was reimplemented. The only numerical comparison remains the inner-loop benchmark against a first-order sliding-mode baseline in Section 3 (Simulation Results). The last cell of the table also states the present limit of the claim: the design uses the quasi-static regime of Assumptionn3, and the evidence in this paper is simulation.
%\textcolor{red}{LINKS!!!!!!!!!!!}
\end{reply}
%\ref{sec:simulations}~\ref{as:smallangle}
\begin{point}
The structure of the adaptive law is not clear enough.
\end{point}

\begin{reply}
We thank the reviewer for this comment. The adaptation law was introduced inside the stability proof, and it was written as the derivative of the parameter error \(\boldsymbol{\tilde{\kappa}}\). That error depends on the unknown matrices \(\boldsymbol{\kappa}_u\) and \(\boldsymbol{\kappa}_v\), so the expression did not show which signals are integrated. The leakage term was attached to that error derivative, with the opposite sign of a restoring \(\sigma\)-modification.

The law is now stated in Section 2.5 (Dynamic Compensator), immediately after the control law (45), as (46). It updates the estimates \(\hat{\boldsymbol{\kappa}}_u\) and \(\hat{\boldsymbol{\kappa}}_v\). Both are \(n\times n\) matrices, with \(n=4\), and the gain \(\gamma>0\) is a scalar. Each equation has two terms. The gradient term is \(\gamma^{-1}\mathbf{u}\mathbf{s}^{T}\) or \(-\gamma^{-1}\mathbf{v}\mathbf{s}^{T}\), and its sign follows the model (6). The leakage term is \(-\|\mathbf{s}\|\hat{\boldsymbol{\kappa}}^{T}\), and it uses only the current estimate. The right-hand side does not contain the unknown matrices.

The proof keeps an equivalent form, (45), obtained from \(\dot{\boldsymbol{\tilde{\kappa}}}=-\dot{\hat{\boldsymbol{\kappa}}}\) because the true parameters are constant. That form is used only to cancel the regressor in \(\dot{V}\). After the substitution, the factor \(\gamma\) multiplies the leakage traces, and the squared parameter errors enter \(\dot{V}\) with a negative sign. \textcolor{red}{Links!!!!!!!!!!!}
\end{reply}

\begin{point}
Why does the regressor need to be persistently exciting for the parameter estimation error to converge to zero?
\end{point}

\begin{reply}
We thank the reviewer for this question. Persistent excitation is required for parameter convergence only in an ideal gradient law, with no disturbance and no leakage. In that law the update is driven by the products of the regressor with the sliding variable. After the sliding variable has reached zero, the update stops. Any component of the parameter error that stays orthogonal to the recorded regressor produces the same closed-loop signals, so it is never corrected. The integral condition
\begin{equation*}
\int_{t}^{t+T}\boldsymbol{\phi}(\tau)\boldsymbol{\phi}(\tau)^{\top}\,d\tau \ge \alpha I
\end{equation*}
removes that ambiguity: the only parameter error compatible with a zero residual is the zero error. For the full matrices \(\boldsymbol{\kappa}_u\) and \(\boldsymbol{\kappa}_v\), the regressor \(\boldsymbol{\phi}\) is the Kronecker regressor built from \(\mathbf{u}\) and \(\mathbf{v}\).
%\eqref{eq:ajuste2}
That conclusion does not transfer to the law of this paper. Equation (57)  keeps a leakage proportional to \(\|\mathbf{s}\|\), and the sliding dynamics keep the disturbance \(\boldsymbol{\delta}(t)\). The leakage balances the gradient away from the true matrices, and the disturbance imitates a parameter error. Both effects remain under persistent excitation. The update is also scaled by signals that vanish with \(\mathbf{s}\), so the estimates freeze in the sliding regime. Theorem 2 therefore states uniform ultimate boundedness of the parameter errors.

The conclusions previously said that parameter convergence required persistent excitation of the regressor. That sentence has been removed. The reasoning above is now Remark 1, placed immediately after the proof of Theorem 2. %~\ref{rem:pe}
\end{reply}

\begin{point}
The authors are suggested to include figures showing error convergence under different gains to illustrate the impact of parameters on tracking performance.
\end{point}

\begin{reply}
The authors thank the reviewer for the positive feedback. We hope that our responses adequately address the questions and suggestions raised.
\end{reply}

\begin{point}
The authors are suggested to provide some recent results on the related topic of adaptive control, such as ‘An equilibrium factor-based iterative learning control of robot arms’, ‘Neural network adaptive iterative learning control for strict-feedback unknown delay systems’.
\end{point}

\begin{reply}
We thank the reviewer for this suggestion. Both works are now cited in the introduction, in a new paragraph that follows the discussion of adaptive and learning-based compensation. They belong to iterative learning control (ILC), an adaptive approach that the previous introduction did not cover.

The paragraph summarizes the two designs. Shen et al. (2025), \emph{IEEE Transactions on Neural Networks and Learning Systems}, 36(7):13460--13469, treat strict-feedback systems with unknown state delays and input saturation. They combine radial-basis-function neural networks, Lyapunov--Krasovskii functionals, command filtering, and auxiliary saturation-compensation systems. Shen et al. (2026), \emph{IEEE Transactions on Reliability}, 75:87--96, introduce update laws with an equilibrium factor that tolerate initial-state errors and actuator faults in robot arms. In both works, convergence is established along the iteration axis through a composite energy function. The same paragraph cites two ILC designs for quadrotors, Schoellig et al. (2012) and Han et al. (2025), which learn over repeated or periodic trajectories.

The paragraph also states how ILC differs from the compensator of this paper. ILC refines the input from one execution of the same task to the next, so it presupposes a task repeated over a finite interval. A delivery flight is executed once, along a reference that changes with each mission. The dynamic compensator therefore adapts its parameters along the time axis, within a single flight. The two suggested works also treat input saturation, delays, and actuator faults explicitly. These effects are not included in our analysis, and the paragraph refers to the conclusions, where their characterization is listed as future work.

While revising this part, we noticed that Han et al. (2025) had been cited among the cooperative multi-UAV transport studies. That work concerns a single quadrotor without a suspended load, so the citation has been moved to the new paragraph.
\end{reply}



\bibliographystyle{IEEEtran}
\bibliography{References}

\end{document}


\begin{point}

\end{point}

\begin{reply}

\end{reply}
